\documentclass[letterpaper, 12pt]{article}

% Use our custom notes template
\usepackage{notes_template}

% --- Document Starts Here ---
\begin{document}

\section*{Updating a Cached QR Factorization with Givens Rotations}

Consider a full-rank matrix $A \in \mathbb{R}^{m \times n}$ with QR factorization
\[
    A = Q R, \qquad R = \begin{pmatrix}
        r_{11} & r_{12} & \cdots & r_{1n} \\
                & r_{22} & \cdots & r_{2n} \\
                &        & \ddots & \vdots \\
                &        &        & r_{nn}
    \end{pmatrix},
\]
and cached right-hand side $c = Q^T b$. Appending a single constraint $a^T x = c_{n+1}$ produces the augmented system
\[
    \begin{pmatrix}
        A \\ a^T
    \end{pmatrix} x' = \begin{pmatrix}
        b \\ c_{n+1}
    \end{pmatrix}.
\]
Multiplying by $Q^T$ keeps the orthogonal factor implicit:
\[
    \begin{pmatrix}
        R \\ w^T
    \end{pmatrix} x' = \begin{pmatrix}
        c \\ c_{n+1}
    \end{pmatrix}, \qquad w^T = a^T Q.
\]
Only the last row breaks triangular form, so we repair it with Givens rotations instead of recomputing the factorization from scratch.

\subsection*{Which rows do we rotate?}

Each rotation acts on a $2\times 2$ block of rows. To eliminate the entry below the diagonal in column $j$, we rotate
\[
    	ext{row } j \text{ and row } n+1,
\]
leaving every other row untouched. The rotation matrix is
\[
    G_{j,n+1}(c,s) = \begin{pmatrix}
        I_{j-1} &              &             &             \\
                 &  c          &             &  s          \\
                 &             & I_{n-j}     &             \\
                 & -s          &             &  c
    \end{pmatrix},
\]
with parameters chosen so that
\[
    \begin{pmatrix}
        c & s \\
        -s & c
    \end{pmatrix}
    \begin{pmatrix}
        r_{jj} \\
        w_j
    \end{pmatrix}
    =
    \begin{pmatrix}
        \hat{r}_{jj} \\
        0
    \end{pmatrix}.
\]
Because $Q$ changes implicitly, we apply the same rotation to the cached right-hand side: $c \leftarrow G_{j,n+1}^T c$.

\subsection*{Column-by-column sweep}

We work from the last column toward the first:
\begin{enumerate}
    \item For $j = n, n-1, \ldots, 1$:
    \begin{enumerate}
        \item Read the pivot $r_{jj}$ from row $j$ and the subdiagonal entry $w_j$ from row $n+1$.
        \item Compute the stable rotation
        \[
            r = \sqrt{r_{jj}^2 + w_j^2}, \qquad c = \frac{r_{jj}}{r}, \qquad s = \frac{w_j}{r}.
        \]
        \item Apply $G_{j,n+1}(c,s)$ to rows $j$ and $n+1$ of $R$ (only columns $j, j+1, \ldots, n$ are touched).
        \item Rotate the matching entries of $w$ and the cached vector $c$.
    \end{enumerate}
\end{enumerate}
After column $j$ is processed, entry $(n+1, j)$ becomes zero. Since each step touches just two rows, the overall update costs $\mathcal{O}(n)$.

\subsection*{Final triangular form}

When the sweep completes we obtain
\[
    	\tilde{R} = \begin{pmatrix}
        r'_{11} & r'_{12} & \cdots & r'_{1n} \\
                 & r'_{22} & \cdots & r'_{2n} \\
                 &         & \ddots & \vdots \\
                 &         &        & r'_{nn} \\
                 &         &        & r_{n+1,n+1}
    \end{pmatrix},
    \qquad
    	\tilde{c} = \begin{pmatrix}
        c'_1 \\
        \vdots \\
        c'_n \\
        c_{n+1}
    \end{pmatrix}.
\]
Back-substitution on the first $n$ rows gives the updated solution $x'$, while the final entry tracks the residual of the new constraint.

\subsection*{Two-by-two picture}

Each step is just the classic $2\times 2$ rotation
\[
    \begin{pmatrix}
        \cos \theta & \sin \theta \\
        -\sin \theta & \cos \theta
    \end{pmatrix}
    \begin{pmatrix}
        a \\
        b
    \end{pmatrix}
    =
    \begin{pmatrix}
        \sqrt{a^2 + b^2} \\
        0
    \end{pmatrix},
\]
applied to rows $j$ and $n+1$. Recognizing that only those two rows move makes the update workflow easy to visualize and implement.

\end{document}
